% called by main.tex
%
\chapter{Conclusiones y Trabajo Futuro}
\label{ch::chapter10}

%----------------------------------------------------------------------
\section{Conclusiones principales}
\label{sec:ch10_conclusiones_principales}
%----------------------------------------------------------------------

Para cerrar el trabajo con las conclusiones, queremos recuperar la pregunta con la que se
inició el planteamiento. Y es que la pregunta que ha guiado este trabajo no era únicamente si un modelo de lenguaje puede
generar texto con apariencia de YAML. Ese sería, en cierto modo, el planteamiento más
superficial del problema. La cuestión de fondo era si un modelo autoregresivo, cuyo
mecanismo natural consiste en producir una secuencia plana de tokens, puede generar una
salida cuya corrección depende de una organización jerárquica: documentos, líneas, listas,
claves, valores y niveles de indentación que deben encajar entre sí. Los resultados del
Capítulo~\ref{ch::chapter9} permiten responder a esa pregunta con una conclusión matizada:
la estructura puede hacerse mucho más fiable cuando se convierte en un objeto explícito de
generación, reconstrucción y evaluación, pero eso no implica que cualquier mecanismo
explícito de predicción jerárquica mejore automáticamente al aprendizaje serializado.

\subsection{Análisis de la propuesta serializada}

El resultado más claro aparece al comparar la línea base con el modelo
\texttt{serialized\_sft}. En test, la tasa de YAML parseable pasa del 46{,}77\,\% al
97{,}14\,\%, y el salto se repite en las métricas que miden contenido, jerarquía y
adecuación al prompt. El F1 de líneas sube de 31{,}75\,\% a 79{,}09\,\%, la exactitud de
\texttt{level} pasa de 22{,}26\,\% a 62{,}98\,\%, el F1 de requisitos del prompt aumenta
de 32{,}30\,\% a 69{,}32\,\% y el F1 de claves semánticas llega al 92{,}11\,\%. Esta
mejora no debe interpretarse solo como una ganancia de formato. Lo importante es que el
modelo aprende a producir una representación intermedia que un parser determinista puede transformar
de forma estable en manifiestos YAML. En otras palabras, la supervisión no solo mejora la
fluidez de la salida, sino que cambia el régimen del sistema: de generar texto
aproximadamente estructurado pasa a generar bloques que conservan suficiente información
jerárquica como para ser reconstruidos. De esta forma, hemos conseguido que un modelo autoregresivo, cuya
naturaleza es secuencial, pueda producir una salida que respeta una jerarquía de niveles y
listas, y que esa jerarquía sea observable, medible y evaluable.

\subsection{Análisis de la propuesta Two-Head}

Esta conclusión da valor a la representación por bloques y al uso de \texttt{level}, pero
no confirma, con esta configuración, la hipótesis de que separar la jerarquía en un cabezal adicional sea,
por sí mismo, la mejor estrategia. El modelo THM v1 mejora a la línea base en varias
dimensiones y obtiene un MAE de nivel competitivo, pero no alcanza la robustez global del
SFT serializado. Su 51{,}61\,\% de YAML parseable queda lejos del 97{,}14\,\% de
\texttt{serialized\_sft}, y la pérdida de parseabilidad arrastra también las métricas de
contenido y adecuación. Esto no invalida la idea de predecir la jerarquía como una señal
estructural diferenciada, pero sí muestra que esa separación introduce una frontera
delicada: el texto de la línea, el nivel predicho y el parser deben quedar suficientemente
acoplados. Si esa coordinación falla, el sistema puede producir fragmentos reconocibles y,
al mismo tiempo, no llegar a un YAML utilizable.

Una primera explicación de este comportamiento está en el desbalanceo de la propia variable
\texttt{level}. Los niveles profundos aparecen con mucha menos frecuencia que los niveles
superficiales, de modo que el cabezal recibe poca señal efectiva para aprender esas clases.
En esas condiciones, es razonable pensar que una pérdida de clasificación estándar favorezca
las regiones más frecuentes de la jerarquía y que el modelo tienda a comprimir la profundidad
hacia los niveles más comunes. Técnicas de reponderación, muestreo o pérdidas más sensibles a la
densidad de cada nivel podrían haber reducido este efecto y mejorado la exactitud del cabezal. Además, los
niveles profundos no son etiquetas raras en abstracto: suelen aparecer precisamente en las
partes más difíciles del dominio, como plantillas de pod, contenedores, comandos, volúmenes,
\texttt{volumeMounts} y listas anidadas. El desbalanceo, por tanto, coincide con una zona de
mayor complejidad estructural.

A esta dificultad se suma el desacoplamiento que introduce la propia arquitectura. En el modelo
\texttt{serialized\_sft}, el nivel forma parte de la misma superficie textual que el contenido
de la línea. El modelo aprende a emitir ambos elementos dentro de una única secuencia
autoregresiva, y puede capturar correlaciones locales entre lo que escribe y la profundidad a
la que suele escribirse. En THM v1, en cambio, el contenido lo genera el cabezal de lenguaje y
el nivel lo predice un cabezal estructural separado. Esa separación es conceptualmente más
limpia, porque convierte \texttt{level} en una variable supervisada explícita, pero también
elimina una dependencia práctica: el nivel predicho no vuelve al generador de texto y, por
tanto, no condiciona directamente las líneas siguientes. De ahí que el modelo pueda producir
contenido Kubernetes reconocible y, aun así, fallar en la reconstrucción final si la trayectoria
de niveles que acompaña a ese contenido no es coherente.

La elección de \texttt{record\_prefix\_state} refuerza esta lectura. En el diseño del THM v1,
el cabezal de nivel lee el estado oculto inmediatamente anterior al registro de una nueva
línea. Esta decisión era metodológicamente atractiva porque evita que el cabezal vea el
contenido de la línea antes de predecir su profundidad. La predicción de \texttt{level} no se
apoya, por tanto, en pistas textuales ya generadas, sino en el estado autoregresivo acumulado
hasta ese punto.

Sin embargo, esa misma limpieza convierte la tarea en una predicción más exigente. En un
manifiesto Kubernetes, muchas transiciones de profundidad dependen de la función concreta de
la línea que empieza: si abre una lista, si continúa un bloque \texttt{command}, si introduce
un \texttt{volumeMount}, si vuelve a un campo hermano o si entra en
\texttt{spec.template.spec}. Al usar \texttt{record\_prefix\_state}, el cabezal debe anticipar
esa decisión antes de observar la línea. Esto puede ser suficiente para niveles frecuentes y
transiciones simples, pero se vuelve frágil en regiones profundas y locales de la jerarquía.

Esta interpretación ayuda también a leer las métricas con más cuidado. La caída global del
THM v1 está muy condicionada por la parseabilidad: cuando una salida no llega a YAML válido,
muchas métricas posteriores quedan arrastradas o directamente dejan de ser informativas. En
los casos en los que la salida sí supera esa frontera, la distancia con \texttt{serialized\_sft}
se reduce. Por tanto, la conclusión no es que la arquitectura sea incapaz de generar contenido
Kubernetes razonable, sino que no consigue mantener de forma estable la interfaz entre línea
generada, nivel predicho y reconstrucción parser-facing. El problema no es solo acertar clases
de \texttt{level} de manera aislada, sino sostener una trayectoria jerárquica completa.

Las variantes ordinales y posicionales refuerzan esta lectura. El objetivo de esos
experimentos era atacar una limitación real de la predicción de niveles, especialmente en
profundidades mayores y en zonas donde el estado latente cambia con la posición de la
línea. Si el estado usado por el cabezal cambia según el momento
del manifiesto, una única frontera global para predecir \texttt{level} puede quedar mal
calibrada. La introducción de codificación posicional y, más tarde, de modulación FiLM,
intentaba precisamente condicionar cómo se interpreta ese estado oculto según la posición de
la línea dentro de la secuencia YAML. Sin embargo, los resultados de THOP-FiLM muestran que
añadir estructura explícita o condicionamiento posicional no basta si la arquitectura completa
no mantiene la estabilidad. En test, esta variante queda por debajo incluso del
THM v1 en las métricas globales.

\subsection{Análisis de la optimización por preferencias}

La fase de DPO ofrece una conclusión distinta. A diferencia de las variantes two-head, DPO
parte del modelo \texttt{serialized\_sft}, es decir, de una política que ya domina
razonablemente el contrato estructural. En ese contexto, las preferencias generadas sí
consiguen mover el modelo en una dirección útil sin romper la superficie de salida. DPO v2
alcanza el 100{,}00\,\% de parseabilidad YAML en test, mejora el F1 de requisitos del
prompt hasta el 72{,}60\,\%, eleva el F1 de claves semánticas hasta el 94{,}50\,\% y
obtiene el mejor \texttt{kubernetes\_domain\_score}, con un 84{,}62\,\%. A pesar de que la
mejora no es absoluta, partimos de un modelo que ya era muy capaz, y el hecho de conseguir
un avance adicional en métricas de contenido y adecuación al prompt es un resultado
significativo. La conclusión es que, una vez que la jerarquía se ha estabilizado, la optimización
de preferencias puede mejorar la adecuación al dominio y la fidelidad a los requisitos del
prompt sin comprometer la estructura.

Este paso fue importante además porque nos permitión obtener un modelo que genera manifiestos Kubernetes
más correctos que los que se encontraban en el propio conjunto de test. Este es un resultado
muy favorable para esta técnica, ya que indica que el modelos se ha aprovechado de conocimiento externo
al conjunto de datos, y que ha aprendido a generalizar más allá de los ejemplos, superando en calidad a las
muestras contra las que se evaluaba.

Por último, el planteamiento iterativo de DPO nos permitió explorar una posible estrategia de mejora online,
manteniendo las ventajas de ligereza y eficiencia de la optimización de preferencias. Pudimos apreciar que
la construcción del conjunto de preferencias en función del aprendizaje de tu política base, y tu objetivo de
entrenamiento, puede ser una estrategia muy efectiva para iterar sobre el modelo y mejorar su desempeño en 
tareas específicas. Aunque por otro lado suponga un coste computacional adicional.

\section{Limitaciones}
Una primera limitación del trabajo está relacionada con el tamaño y la diversidad del conjunto de entrenamiento.
Aunque el conjunto utilizado permite estudiar el problema de la generación estructurada de manifiestos Kubernetes de
forma controlada, su cobertura del dominio sigue siendo limitada. Kubernetes admite múltiples tipos de recursos,
combinaciones entre ellos, campos opcionales, convenciones de despliegue y patrones de configuración que no aparecen
con la misma frecuencia en los datos disponibles. Esta limitación afecta directamente al alcance del aprendizaje: el
modelo puede captar regularidades importantes del formato YAML y de la jerarquía de los manifiestos, pero trabaja sobre
una representación parcial de la variedad real del ecosistema Kubernetes.

Una segunda limitación aparece en las métricas específicas de Kubernetes obtenidas sobre el propio conjunto de entrenamiento.
Este punto es especialmente relevante porque desplaza el problema más allá de la generalización a ejemplos no vistos.
Si determinadas métricas del dominio siguen siendo bajas incluso en los datos empleados para el ajuste, la dificultad no
está solo en extrapolar, sino también en representar correctamente las relaciones internas del manifiesto: nombres de
recursos, campos esperados, dependencias entre secciones y coherencia semántica básica. En este sentido, los resultados
refuerzan una de las ideas centrales del trabajo: generar una salida con apariencia de YAML no equivale necesariamente
a producir un manifiesto Kubernetes correcto desde el punto de vista estructural y funcional.

La capacidad de cómputo disponible también condicionó el desarrollo experimental. El entrenamiento y la evaluación de
modelos de lenguaje, incluso con técnicas de adaptación eficiente como LoRA o con modelos cuantizados, requieren
recursos significativos de memoria, tiempo y almacenamiento. En la práctica, esta restricción redujo el número de
configuraciones exploradas, la repetición de experimentos con distintas semillas, la búsqueda de hiperparámetros y
la comparación sistemática entre modelos base de distinto tamaño. Por tanto, el trabajo se centra en una ruta
experimental concreta, suficiente para analizar la comparación entre las formulaciones propuestas, pero todavía
limitada frente a una exploración más amplia del espacio de entrenamiento.

Por último, la anotación manual de preferencias introduce una limitación de tiempo y de criterio experto. El proceso
de preferencia humana ha permitido valorar mejor aspectos como la adecuación al prompt, la coherencia semántica del
manifiesto o la utilidad práctica de la configuración generada. Sin embargo, este tipo de anotación exige revisar
salidas técnicas con detalle y aplicar criterios consistentes sobre un dominio que no siempre admite una decisión trivial.
Esta dificultad explica que el trabajo se apoye en señales semiautomáticas de preferencia para las etapas
posteriores de optimización, donde he necesitado el uso de métricas de apoyo para generar preferencias de forma más eficiente y escalable.

\section{Revisión de objetivos}



\section{Trabajo futuro}
